\statement{A.129}{Proved}{The actual boundary norm, complete reflection cohomology and the specified cusp transporter}
Retain the original negative full-GT data and canonical word.
The boundary defect \(D=[\omega,\xi]p^{-d}\) is an explicitly
specified derived cocycle for a noncusp reflection whose entire
derived \(H^1\) is trivial. The actual norm-preserving correction
subgroup is nevertheless trivial, and the norm word is unique.
The cusp actions have instead the whole pointed compact cohomology
space \(\widehat{\mathbb Z}\). With the original C.13 marking, the
boundary remains a specified conjugacy; equivalently it is equality
of the printed canonical word with the unique noncusp fixed generator
supplied by A.130. No boundary, full-generation or original CM
vanishing is concluded.
\begin{proof}
\emph{A.129.1. Actual data and conventions}

Let E=Fhat(X,Y), let E' be its closed derived subgroup, and put
\[
 p=[Y,X],\quad q=X,\quad z=(pq)^{-1}=YX^{-1}Y^{-1}.
\]
All commutators mean [a,b]=\(aba^{-1}\)\(b^{-1}\). The original negative full-GT involution has parameter phi, canonical commutator word u, actual even k, and actual full unit d=2k+1. Retain
\[
 U=u(X,p),\quad V=u(z,p),\quad\Phi=\phi(X,p),\quad Q=\phi(z,p).
\]
The actual elliptic section involution Theta satisfies
\[
 \Theta X=\Phi X^{-1}\Phi^{-1},\qquad
 \Theta Y=p^{-1}QY\Phi^{-1},\qquad
 \Theta p=p^{-1},\qquad\Theta^2=1.
\tag{A.129.1}
\]
The original norm and duality, with the marked pair embeddings retained, give
\[
 \Theta U=U\Phi^{-1},\qquad
 \Theta V=p^{-1}VQ^{-1}p.
\tag{A.129.2}
\]
These are the complete inputs already proved in the original cusp bridge; no section is being assigned to the unproved pair (d,u).

Write
\[
 P=p^d,\qquad A=\xi=UX^dU^{-1},\qquad
 W=\omega=p^kVYU^{-1},\qquad D=[W,A]P^{-1}.
\tag{A.129.3}
\]
Here P is an element of E, not the pure elliptic braid group. From (A.129.1)--(A.129.2), in the original order,
\[
 \Theta A=A^{-1},\qquad\Theta W=P^{-1}W,
 \qquad\Theta P=P^{-1}.
\tag{A.129.4}
\]
The requested boundary equation is D=1. It is equivalent to the specified original third-cusp conjugacy printed in the earlier bridge; it is not just an unspecified peripheral-conjugacy condition.

Define the actual involution
\[
 \beta_2=\operatorname{Ad}(\Phi^{-1})\Theta.
\]
Since Theta(Phi)=\(\Phi^{-1}\), it has order two, and its restriction to the embedded original pair (q,p) is
\[
 \beta_2(q)=q^{-1},\qquad
 \beta_2(p)=\Phi^{-1}p^{-1}\Phi.
\tag{A.129.5}
\]
We never replace beta2 by simultaneous inversion theta or by standard elliptic conjugation.

\par\smallskip\noindent\emph{A.129.2. The exact boundary coboundary}

Put
\[
 g=PAP,\qquad t=PW^{-1},\qquad
 \kappa_D=\operatorname{Ad}(g)\Theta,
 \qquad\kappa_P=\operatorname{Ad}(P)\Theta.
\tag{A.129.6}
\]
Equation (A.129.4) gives Theta(g)=\(g^{-1}\); both maps in (A.129.6) are involutions. Since C=[W,A]=DP,
\[
\begin{aligned}
 \Theta D
 &=[P^{-1}W,A^{-1}]P\\
 &=P^{-1}A^{-1}C^{-1}PAP\\
 &=(PAP)^{-1}D^{-1}(PAP).
\end{aligned}
\tag{A.129.7}
\]
The second equality uses W \(A^{-1}\)\(W^{-1}\)=\(A^{-1}\)\(C^{-1}\); the third uses \(C^{-1}\)=\(P^{-1}\)\(D^{-1}\). In particular D is a kappaD-cocycle.

There is the stronger literal identity
\[
 \boxed{D=t^{-1}\kappa_D(t).}
\tag{A.129.8}
\]
Indeed Theta(t)=\(P^{-1}\)\(W^{-1}\)P, and substitution in (A.129.8) cancels to W A \(W^{-1}\) \(A^{-1}\)\(P^{-1}\). This is a coboundary in E. Its values belong to E', while
\[
 [t]=-[Y]\in E_{\rm ab}=\widehat{\mathbb Z}[X]\oplus\widehat{\mathbb Z}[Y].
\]
Thus the derived cohomology class [D] is the connecting image of -[Y] for the actual action kappaD.

We use \(Z^1(C_2,G;\kappa)=\{z:z\kappa(z)=1\}\), with equivalence \(z\sim b^{-1}z\kappa(b)\). Its neutral cocycles are \(h^{-1}\kappa(h)\).

\par\smallskip\noindent\emph{A.129.3. Neutrality, a specified fixed word, and the restricted repair}

Before invoking C.8/C.11, the exact general connecting-map criterion is
\[
 [D]=1\text{ in }H^1(C_2,E';\kappa_D)
 \iff\exists h\in E^{\kappa_D}\text{ with }[h]=-[Y].
\tag{A.129.9}
\]
For example if D=\(b^{-1}\)kappaD(b) with b in E', then t \(b^{-1}\) is fixed and has the required homology. Conversely a fixed h with the same homology gives \(t^{-1}\)h in E' and kills the cocycle by that change of lift. All fixed elements have zero X-coordinate because kappaD acts on abelianization by diag(-1,1). Thus (A.129.9) is equivalently surjectivity of the fixed subgroup onto the full Y-line. It is not the original boundary equation, which is precisely
\[
 D=1\iff\kappa_D(t)=t
\tag{A.129.10}
\]
for the specified t.

A proposed change t\(\to\)tb, with b in E', corresponds exactly to W\(\to\)W'=\(b^{-1}\)W. Preserving the actual norm equation in (A.129.4) requires
\[
 \Theta W'=P^{-1}W'
 \iff\Theta b=P^{-1}bP
 \iff b\in(E')^{\kappa_P}.
\tag{A.129.11}
\]
Under this additional condition, the actual new boundary defect is
\[
 D'=b^{-1}D\kappa_D(b).
\]
Since kappaD=Ad(PA)kappaP, the constrained repair equation is
\[
 \exists b\in(E')^{\kappa_P}:D=[b,PA].
\tag{A.129.12}
\]
The norm identity also gives
\[
 \kappa_P=\operatorname{Ad}(W)\Theta\operatorname{Ad}(W^{-1}),
\tag{A.129.13}
\]
so permitted corrections have the form b=\(WcW^{-1}\) with c in \((E')^\Theta\), and W'=\(Wc^{-1}\). No such correction is inferred from ambient or derived cohomology alone.

\par\smallskip\noindent\emph{A.129.4. Fresh full scope audit of C.8, C.11, and their inputs}

The exact previously proved inputs are C.5--C.8 and C.11--C.13, at their full profinite scopes.

C.8 assumes an outer involution of the full free profinite punctured-torus group preserving the commutator cusp with FULL multiplier -1. It constructs the cusp-normalized extension, distinguishes one versus two section classes by the dyadic homology action, fuses the two local cusp reflections, and proves that the cusp section has trivial geometric centralizer. Its proof uses the actual peripheral duality in C.6 and the finite permutation-module fixed-coset witness of C.7; it makes no assumption that the automorphism comes from a discrete surface homeomorphism.

The actual Theta satisfies these hypotheses literally: its outer class has order two because its abelian action is diag(-1,1), and it sends p=[Y,X] to \(p^{-1}\). Interchanging the names of the free generators matches the printed cusp [x,y]. Its dyadic matrix is I modulo two. Hence the actual extension E semidirect <Theta> has exactly two E-conjugacy classes of sections. The cusp section Theta has
\[
 \boxed{E^\Theta=1.}
\tag{A.129.14}
\]
The other section has a geometric centralizer B whose homology image, by the ENTIRE C.11 proof, is the full fixed line Zhat[Y]. This is stronger than a nonzero mod-two image: odd cyclic covers establish surjectivity on all odd invariant quotients; the nilpotent localized class establishes a dyadic unit; the closed image is then the full line.

The anti-invariant character X modulo two distinguishes the two section classes. Cusp conjugates have coefficient \((I-\Theta_{\mathrm{ab}})[h]=2h_X[X]\) in abelianization, so they have even X-label. The section gTheta in (A.129.6) has coefficient [g]=d[X], with d odd. It is therefore the noncusp class.

\par\smallskip\noindent\emph{A.129.5. Complete derived H1: both actual branches}

For the cusp action Theta, every E'-valued cocycle gives a section with even X-label, hence belongs to the neutral ambient section class. The exact connecting map from the fixed homology line therefore surjects onto \(H^1\)(C2,E';Theta). It is injective because \(E^\Theta\)=1. Thus
\[
 \boxed{H^1(C_2,E';\Theta)\cong\widehat{\mathbb Z}[Y]}
\tag{A.129.15}
\]
as pointed compact spaces.

For clarity, the topology does not presume a continuous choice of lifts. The compact set of lifts in E with homology on the fixed line maps continuously to the compact cocycle quotient; two lifts with the same homology differ in E' and give equivalent cocycles. The map descends through the compact quotient to the line. The inverse is continuous because a continuous bijection from a compact space to the Hausdorff quotient by a compact group action is a homeomorphism.

Equation (A.129.2) gives the exact conjugacy
\[
 \beta_2=\operatorname{Ad}(U^{-1})\Theta\operatorname{Ad}(U),
 \qquad U\in E'.
\tag{A.129.16}
\]
Consequently \(E^{\beta_2}\)=1 and the same full pointed compact space (A.129.15) describes \(H^1\)(C2,E';beta2), with the fixed homology marking unchanged.

For the NONCUSP action kappaD, let z be any E'-valued cocycle. The section z(gTheta) still has odd X-label. There are exactly two classes, so it is E-conjugate to gTheta and z=\(h^{-1}\)kappaD(h) for some h in E. Abelianization forces \(h_X\)=0. C.11 provides q in \(E^{\kappa_D}\) with [q]=[h]. Then h'=\(q^{-1}\)h belongs to E' and
\[
 h'^{-1}\kappa_D(h')=z.
\]
This proves the ENTIRE result
\[
 \boxed{H^1(C_2,E';\kappa_D)=*.}
\tag{A.129.17}
\]
In particular [D] is automatically neutral. The fixed subgroup \(E^{\kappa_D}\) is not asserted trivial; it surjects onto the full Y-line. Formula (A.129.17) is therefore fully consistent with (A.129.14).

Combining (A.129.13)--(A.129.14) gives \(E^{\kappa_P}\)=1. Thus the permitted correction subgroup in (A.129.11)--(A.129.12) is trivial: b=1 is the only norm-preserving repair, and W'=W. More generally the equation \(\Theta(W)=P^{-1}W\) has at most one solution in all E, since the ratio \(W_1^{-1}W_2\) of any two solutions is Theta-fixed. Its actual solution is (A.129.3).

\par\smallskip\noindent\emph{A.129.6. Exact C.13 marking and the remaining specified conjugacy}

C.8 supplies the unique \(g_p\) with
\[
 g_p\Theta g_p^{-1}=p\Theta,
 \quad\text{equivalently }\Theta(g_p)=p^{-1}g_p.
\]
The word \(p^k\) \(g_p\) satisfies the norm equation for P=\(p^{2k+1}\); uniqueness proves
\[
 \boxed{W=p^kg_p,\qquad g_p=VYU^{-1}.}
\tag{A.129.18}
\]
Its abelianization is literally [Y], consistent with C.11's primitive-transporter theorem.

C.13's actual notation is c=p, v=Theta, mu=d and i=0. Its transporter \(a_{d,0}\)=\(p^{(d-1)/2}\)\(g_p\) is therefore exactly W. The boundary equation has the equivalent form
\[
 \boxed{D=1\iff \xi W\xi^{-1}=p^{-d}W=\Theta W.}
\tag{A.129.19}
\]
Thus it is C.13.2 with the additional SPECIFIED conjugator xi. C.13 does not assert even the unspecified-conjugator existence; C.14 retains that issue as Open. Formula (A.129.19) is a correct comparison to that earlier gap, not a way around it. Equivalently the specified outside involution \(\xi^{-1}\Theta\) would centralize W.

The action kappaD is also explicitly related to the actual beta2. In the extension by Theta,
\[
 \xi\Theta=U X^d\beta_2 U^{-1},\qquad
 X^d\beta_2=X^k(X\beta_2)X^{-k}.
\]
With h=\(PUX^k\) this gives
\[
 \kappa_D=\operatorname{Ad}(h)
              (\operatorname{Ad}(X)\beta_2)
              \operatorname{Ad}(h^{-1}),
\]
and the specified test element becomes
\[
 h^{-1}th=X^{-k}Y^{-1}V^{-1}p^{k+1}UX^k.
\tag{A.129.20}
\]
No standard-reflection replacement has been made.

\par\smallskip\noindent\emph{A.129.7. The actual Brunnian corrected class has the OTHER branch}

Retain the completely proved elliptic Brunnian notation: j is the based half-braid action, theta=Ad(\(Y^{-1}\))j is simultaneous inversion in the basis (X,\(YX^{-1}\)), beta2 commutes with theta, and h1 is the unique derived transporter from Theta(Pi0) to Pi0. Its three normalized forgetting images determine b0 in E', with beta2(b0)b0=1. The actual canonical transporter \(p_u\) has common residual word \(a_u\) in E'.

The corrected action and actual C3 parameter are
\[
 \kappa_0=\operatorname{Ad}(b_0^{-1})\beta_2,
 \qquad\ell=b_0^{-1}\beta_2(a_u^{-1})a_u.
\tag{A.129.21}
\]
The previously proved complete calculation gives
\[
 \kappa_0(\ell)\ell=1,\qquad
 \ell=\kappa_0(a_u)^{-1}b_0^{-1}a_u,
 \qquad[\ell]=[b_0^{-1}]
 \text{ in }H^1(C_2,E';\kappa_0).
\tag{A.129.22}
\]
The section for kappa0 in the Theta extension is \(b_0^{-1}\Phi^{-1}\Theta\). Its anti-invariant mod-two label is zero. It is therefore the cusp class, so
\[
 E^{\kappa_0}=1,\qquad
 H^1(C_2,E';\kappa_0)\cong\widehat{\mathbb Z}[Y].
\tag{A.129.23}
\]
Its specified class \([b_0^{-1}]\) has not been evaluated. In particular (A.129.17) does not neutralize (A.129.22): the twisting classes have different abelian mod-two labels. The statement remains distinct even though all these actions have the same linear action on \(E_{\mathrm{ab}}\).

\par\smallskip\noindent\emph{A.129.8. What the exact commutator theorem does and does not constrain}

The actual Brunnian joint kernel now has the complete description Br=closed[U0,V0] in the two-retraction fiber (here U0,V0 denote the actual two closed kernels, not the canonical word U). Independently, the entire connecting space is
\[
 \partial:E'\xrightarrow{\sim}H^1(C_3,\mathrm{Br}).
\]
Thus the commutator presentation is compatible with every C3 parameter in E'. It does not alone imply ell=1.

There is a precise all-word test of the reduced constraints. Fix the actual
\(U,b_0,\beta_2,\theta,\Pi_0,h_1\). For arbitrary \(a\in E'\),
surjectivity of the actual joint forgetting map \(P'\to(E')^3\)
supplies \(p\) with tuple \((Ua,U\theta(a),Ua)\). Define
\(\Pi=p\Pi_0p^{-1}\) and \(h=\Theta(p)h_1p^{-1}\). Then
\(\Pi^3=1\), \(\Theta(\Pi)=h\Pi h^{-1}\), \(\Theta(h)h=1\),
and all previously derived tuple calculations hold with
\(t_a=\beta_2(a)b_0a^{-1}\). The cyclic defect lies in
\(\mathcal B\) and has parameter \(b_0^{-1}\beta_2(a^{-1})a\).
Every identity just listed follows from the actual maps and the chosen
\(p\), not from a finite approximation.

This is ONLY a family satisfying the reduced cubic, tuple and norm identities. The constructed Pi is not asserted to equal the canonical explicit word \(\Pi_u\); no profinite word having the original pentagon is reconstructed from a. It is therefore not an actual GT counterexample or a proof that the full pentagon cannot constrain \(a_u\). The remaining gate must use additional information tying the specified canonical word to its unique transporter, or an exact comparison with the specified boundary conjugacy (A.129.19).

No direct forgetting homomorphism can evaluate a nontrivial boundary defect from a Br element: all three maps kill Br by definition. The parameter ell is extracted using the connecting lift outside Br. The full commutator presentation must be used with that lift and the genuine cyclic action, which does not preserve the chosen third-forgetting fiber in general.

\par\smallskip\noindent\emph{A.129.9. The actual fixed generator.}
The subsequent independent A.130 proves, already for the original
negative \(\widehat{GT}_0\) two-strand action, that
\[
 b=\operatorname{Ad}(X)\beta_2,\qquad
 t_*=\operatorname{Ad}(X)\theta,\qquad
 B=E^b=\overline{\langle a_B\rangle}\cong\widehat{\mathbb Z},
 \qquad [a_B]=[Y].
\]
Here \(\theta=\operatorname{Ad}(Y^{-1})j\) is the actual normalized
half-twist; it commutes with \(\beta_2\). The generator \(a_B\) of
Y-homology one is unique, and \(t_*(a_B)=a_B^{-1}\). This use of
A.130 is confined to the actual two-strand action: no three-strand
\(\widehat{GT}_0\) action or pentagon is inferred. In the following
formulas write \(\beta=\beta_2\), \(b=\operatorname{Ad}(X)\beta\)
and retain the distinct notation \(t_*\) for that involution, separate
from the boundary test element \(t=PW^{-1}\) in paragraph2.
\par\smallskip\noindent\emph{A.129.10. The literal all-exponent family}

For every s in Zhat define
\[
 a(s)=X^{s+1}U^{-1}p^sVY X^{-s-1}.
\tag{A.129.24}
\]
It is a continuous family and [a(s)]=[Y]. The following identities hold for EVERY s:
\[
 \boxed{b(a(s))=a(-s-1),\qquad
        t_*(a(s))=a(-s-1)^{-1}.}
\tag{A.129.25}
\]
No substitution of a new canonical word or a finite parameter scan is involved.

For the first, put \(Q=\phi(YX^{-1}Y^{-1},p)\). The actual norm formulas give
\[
 \beta(U)=\Phi^{-1}U,\quad
 \beta(p)=\Phi^{-1}p^{-1}\Phi,\quad
 \beta(V)=\Phi^{-1}p^{-1}VQ^{-1}p\Phi,\quad
 \beta(Y)=\Phi^{-1}p^{-1}QY.
\]
Consequently
\[
 \beta(a(s))=X^{-s-1}U^{-1}p^{-s-1}VY X^{s+1}.
\]
Conjugating by X gives the first formula in (A.129.25).

For the second, \(j(U)=V\), \(j(V)=pUp^{-1}\), \(\theta=\operatorname{Ad}(Y^{-1})j\) and \(\theta(Y)=XY^{-1}X^{-1}\) give
\[
\begin{aligned}
 t_*(U)&=XY^{-1}VYX^{-1},\\
 t_*(V)&=XY^{-1}pUp^{-1}YX^{-1},\\
 t_*(p)&=XY^{-1}pYX^{-1},\\
 t_*(Y)&=X^2Y^{-1}X^{-2}.
\end{aligned}
\]
Substituting them in (A.129.24) and using \(p^{-1}YXY^{-1}=X\) yields
\[
 t_*(a(s))=X^{-s}Y^{-1}V^{-1}p^{s+1}UX^s,
\]
which is exactly \(a(-s-1)^{-1}\). All profinite powers are images of continuous procyclic homomorphisms; the cancellations are literal and do not reorder powers of different elements.

The composite \(bt_*=\beta\theta\) therefore satisfies
\[
 (\beta\theta)(a(s))=a(s)^{-1}
\tag{A.129.26}
\]
for every s. This is a real additional cocycle identity for the correct composite action, but (A.129.26) does not make a(s) fixed by b.

\par\smallskip\noindent\emph{A.129.11. Equality with the unique fixed generator}

Use the exact boundary notation
\[
 P=p^d,\qquad A=UX^dU^{-1},\qquad
 W=p^kVYU^{-1},\qquad D=[W,A]P^{-1}.
\]
The canonical norm gives \(\Theta(W)=P^{-1}W\). Direct rearrangement shows
\[
 D=1\iff AWA^{-1}=\Theta(W)
      \iff (\operatorname{Ad}(A^{-1})\Theta)(W)=W.
\tag{A.129.27}
\]
In the actual extension by Theta, let \(h_-=UX^{-k-1}\). Since d=2k+1,
\[
 A^{-1}\Theta
 =UX^{-d}\beta U^{-1}
 =h_-(X\beta)h_-^{-1}.
\tag{A.129.28}
\]
Moreover
\[
 h_-^{-1}Wh_-=X^{k+1}U^{-1}p^kVY X^{-k-1}=a(k).
\]
Thus (A.129.27)--(A.129.28), followed by the unique Y-homology generator of B, prove the exact equivalences
\[
\boxed{\begin{aligned}
 D=1
 &\iff a(k)\in B\\
 &\iff a(k)=a_B\\
 &\iff a(k)=a(-k-1)\\
 &\iff t_*(a(k))=a(k)^{-1}.
\end{aligned}}
\tag{A.129.29}
\]
The third and fourth lines use (A.129.25). These are SPECIFIED word equalities. Equal Y-homology does not imply them. Neither the procyclicity of B nor the fact that (A.129.26) already holds puts the specified word a(k) in B.

The earlier full result \(H^1(C_2,E';\kappa_D)=*\) is compatible with (A.129.29): it concerns unrestricted derived cobounders for a different based reflection. The original norm-preserving repair subgroup is trivial. No change of the canonical W, u, k or the original CM frame is authorized by these equivalences.

\end{proof}

\paragraph{Source record and exact remaining scope.}
The complete retained C.8/C.11/C.13 proofs apply to the actual full
commutator cusp and the actual homology actions. The original based
bridge is A.122; A.127 retains the distinct three-strand parameter.
The fixed-generator input is the separately proved A.130, not an
assumption of generation. All displayed word calculations are
symbolic and continuous in their profinite exponents. The original
boundary and actual CM parity remain Open.
